\documentclass[a4paper, serif, 10pt]{moderncv}

%% moderncv
% color and style
\moderncvcolor{black}
\moderncvstyle{banking}

%% page layout/margins
\usepackage[top=2.5cm, bottom=2.5cm, left=1.5cm, right=1.5cm]{geometry}

\nopagenumbers{}

%% Babel config for French language
% ref:
% - https://latex3.github.io/babel/guides/locale-french.html
\usepackage[french]{babel}

%% fontspec
\usepackage{fontspec}

\IfFontExistsTF{Linux Libertine}{
  \setmainfont[Ligatures={TeX, Common}, Numbers=OldStyle]{Linux Libertine}
}{}
\IfFontExistsTF{Linux Libertine O}{
  \setmainfont[Ligatures={TeX, Common}, Numbers=OldStyle]{Linux Libertine O}
}{}

%% CTeX
% CJK support, used to show CJK characters in the resume
%
% - heading=false: disable ctex heading
% - punct=kaiming: use kaiming punctuations styles for CJK
% - scheme=plain: use plain scheme, do not override \normalsize font size
% - space=auto: space settings for CJK characters
%
% CJK fonts are resolved per locale language, e.g. Noto Serif/Sans CJK JP
% for Japanese resumes, with fallback to the next region if not installed
%
% ref:
% - http://ctan.mirrorcatalogs.com/language/chinese/ctex/ctex.pdf
\usepackage[UTF8, heading=false, punct=kaiming, scheme=plain, space=auto]{ctex}

\IfFontExistsTF{Noto Serif CJK SC}{
  \setCJKmainfont{Noto Serif CJK SC}
}{}
\IfFontExistsTF{Noto Sans CJK SC}{
  \setCJKsansfont{Noto Sans CJK SC}
}{}

%% line spacing
\usepackage{setspace}
\setstretch{1.125}

%% URL styling - use normal text instead of monospace
\urlstyle{same}

%% Patch \httplink and \httpslink to handle full URLs
% The original moderncv macros always prepend http:// or https:// to URLs.
% This causes issues when the URL already has a protocol (e.g., https://example.com)
% resulting in malformed URLs like https://https://example.com.
% This patch checks if the URL already contains "://" and uses it directly if so.
% Note: We use \str_set:Nx (x-expansion) to expand macros like \@homepage first.
\ExplSyntaxOn
\str_new:N \g_yamlresume_url_str

\RenewDocumentCommand{\httpslink}{O{}m}{
  \str_set:Nx \g_yamlresume_url_str {#2}
  \str_if_in:NnTF \g_yamlresume_url_str {://}
    { \url{#2} }
    { \url{https://#2} }
}

\RenewDocumentCommand{\httplink}{O{}m}{
  \str_set:Nx \g_yamlresume_url_str {#2}
  \str_if_in:NnTF \g_yamlresume_url_str {://}
    { \url{#2} }
    { \url{http://#2} }
}
\ExplSyntaxOff

\name{Camille Lefèvre}{}
\title{Infirmière diplômée d'État}
\phone[mobile]{+33 6 12 34 56 78}
\email{camille.lefevre@example.fr}
\homepage{https://www.camille-lefevre-infirmiere.fr}

\address{45 Rue de la République, Lyon, Auvergne-Rhône-Alpes, France, 69003}{}{}

\extrainfo{{\small \faLinkedin}\ \href{https://www.linkedin.com/in/camille-lefevre}{@camille-lefevre} {} {} {} • {} {} {} 
{\small \faInstagram}\ \href{https://www.instagram.com/camille.infirmiere}{@camille.infirmiere}}

\begin{document}

\maketitle

\section{Informations de base}

\cvline{}{\begin{itemize}
\item Infirmière diplômée d'État avec 5 ans d'expérience en milieu hospitalier, spécialisée en soins généraux et urgences
\item Excellente capacité à travailler en équipe pluridisciplinaire et à communiquer avec les patients et leurs familles
\item Maîtrise des protocoles de soins, de l'hygiène hospitalière et de la gestion des dossiers patients informatisés
\item Reconnue pour son sens de l'écoute, sa rigueur et sa réactivité face aux situations d'urgence
\end{itemize}}
\section{Formation}

\cventry{sept. 2017–juil. 2020}
        {Licence, Soins Infirmiers, Score : Mention Bien}
        {Institut de Formation en Soins Infirmiers (IFSI) de Lyon}
        {\url{https://www.ifsi-lyon.fr}}
        {}
        {\begin{itemize}
\item Formation complète aux soins infirmiers, incluant stages en hôpital, EHPAD et soins à domicile
\item Acquisition de compétences techniques et relationnelles pour la prise en charge globale du patient
\item Participation à des projets de recherche en soins infirmiers
\end{itemize}
\textbf{Cours} : Sciences infirmières, Anatomie et physiologie, Pharmacologie, Soins d'urgence, Hygiène hospitalière, Relation d'aide, Santé publique, Stages cliniques}
\section{Expérience professionnelle}

\cventry{mars 2023–Present}
        {Infirmière diplômée d'État}
        {CHU de Lyon - Hôpital Édouard Herriot}
        {\url{https://www.chu-lyon.fr}}
        {}
        {\begin{itemize}
\item Prise en charge de patients en service de médecine polyvalente (30 lits)
\item Réalisation des soins techniques : prises de sang, injections, pansements, sondages
\item Gestion des dossiers patients informatisés et coordination avec les médecins et aides-soignants
\item Éducation thérapeutique des patients et de leurs familles
\item Participation à la démarche qualité et à la lutte contre les infections nosocomiales
\end{itemize}
\textbf{Mots-clés} : Soins généraux, Médecine polyvalente, Éducation thérapeutique, Dossier patient informatisé}

\cventry{sept. 2020–févr. 2023}
        {Infirmière diplômée d'État}
        {Clinique de la Sauvegarde}
        {\url{https://www.clinique-sauvegarde.fr}}
        {}
        {\begin{itemize}
\item Soins infirmiers en service de chirurgie orthopédique (25 lits)
\item Préparation et suivi des patients avant et après opération
\item Administration des traitements antalgiques et surveillance post-opératoire
\item Collaboration avec les chirurgiens, kinésithérapeutes et équipes de rééducation
\item Accompagnement des patients dans leur parcours de soins et de rééducation
\end{itemize}
\textbf{Mots-clés} : Chirurgie orthopédique, Soins post-opératoires, Gestion de la douleur, Rééducation}
\section{Langues}

\cvline{Français}{Niveau natif ou bilingue}
\cvline{Anglais}{Niveau professionnel complet \hfill \textbf{Mots-clés} : TOEIC 850, Communication médicale}
\cvline{Espagnol}{Niveau élémentaire}
\section{Compétences}

\cvline{Soins infirmiers}{Expert \hfill \textbf{Mots-clés} : Prise de constantes, Injections, Pansements, Sondages, Prélèvements sanguins}
\cvline{Urgences}{Avancé \hfill \textbf{Mots-clés} : Gestion du stress, Protocoles d'urgence, AFGSU 2}
\cvline{Relation d'aide}{Avancé \hfill \textbf{Mots-clés} : Écoute active, Communication, Éducation thérapeutique}
\cvline{Informatique médicale}{Intermédiaire \hfill \textbf{Mots-clés} : DPI, Logiciel Hôpital, Microsoft Office}
\section{Récompenses}

\cventry{juil. 2020}
        {IFSI de Lyon}
        {Prix d'excellence en soins infirmiers}
        {}
        {}
        {Récompense décernée pour la qualité des soins prodigués lors des stages et l'implication dans la formation.}
\section{Certificats}

\cventry{sept. 2021}
        {CESU de Lyon}
        {AFGSU 2 - Attestation de Formation aux Gestes et Soins d'Urgence}
        {\url{https://www.cesu-lyon.fr}}
        {}
        {}

\cventry{janv. 2022}
        {Hospices Civils de Lyon}
        {Formation aux bonnes pratiques d'hygiène}
        {\url{https://www.chu-lyon.fr/formation-hygiene}}
        {}
        {}

\cventry{juin 2023}
        {Université Claude Bernard Lyon 1}
        {Certification en éducation thérapeutique du patient}
        {\url{https://www.univ-lyon1.fr}}
        {}
        {}
\section{Publications}

\cventry{oct. 2022}
        {Prise en charge infirmière de la douleur post-opératoire : une revue de la littérature}
        {Congrès National des Infirmiers}
        {\url{https://www.congres-infirmiers.fr}}
        {}
        {\begin{itemize}
\item Poster présenté lors du congrès annuel
\item Analyse des pratiques et recommandations pour l'amélioration de la prise en charge
\end{itemize}}
\section{Projets}

\cventry{janv. 2023–juin 2023}
        {Mise en place d'un protocole d'évaluation systématique de la douleur en service de chirurgie orthopédique}
        {Projet d'amélioration de la prise en charge des patients douloureux}
        {\url{https://www.chu-lyon.fr/projet-douleur}}
        {}
        {\begin{itemize}
\item Élaboration d'un outil d'évaluation de la douleur (EVA) et formation des équipes
\item Réduction de 30\% des douleurs non contrôlées signalées
\end{itemize}
\textbf{Mots-clés} : Douleur, Qualité, Protocole}
\section{Centres d'intérêt}

\cvline{Sports}{Randonnée, Yoga, Natation}
\cvline{Culture}{Lecture, Cinéma, Musique}
\section{Bénévolat}

\cventry{janv. 2021–Present}
        {Infirmière bénévole}
        {Croix-Rouge française - Unité locale de Lyon}
        {\url{https://www.croix-rouge.fr}}
        {}
        {\begin{itemize}
\item Participation à des maraudes et à des consultations infirmières pour personnes précaires
\item Réalisation de soins de base, prévention et orientation vers les structures de soins
\item Sensibilisation aux gestes de premiers secours
\end{itemize}}
    
\end{document}